% GENERATED FILE. Edit canonical lessons, then run npm run generate:books.
% canonical-source-hash: fnv1a64:81608b024b4de588
% canonical-lessons: GE-C391-steuer, GE-C391-zins, GE-C391-investieren, GE-C391-anschaffen, GE-C391-versichern

\chapter{Steuern und Zinsen}
\label{ch:ge-steuern-und-zinsen}
\hlchaptermodality{391}

\begin{chapteropening}
\textbf{By the end of this chapter:} \emph{I can talk about taxes and interest, investing, buying something big and taking out insurance.}

Complete the last lesson of chapter 391: I can talk about taxes and interest, investing, buying something big and taking out insurance.
\end{chapteropening}

\section[die Steuer]{die Steuer — the tax}

\label{lesson:GE-C391-steuer}



\begin{quote}
Before the new one: say the German for the fortune, then the German for the payment.
\end{quote}

\subsection*{You'll want to know: die Steuer}
\textbf{die Steuer} — \textquotedblleft{}the tax\textquotedblright{}.

\begin{quote}
Auf Zigaretten zahlt man hohe Steuern.
\end{quote}

\begin{quote}
\emph{You pay high taxes on cigarettes.}
\end{quote}

Plural: \emph{die Steuern}. Careful: \emph{das Steuer} is the steering wheel.

\begin{cousinweb}
\emph{Steuer} \textquotedblleft{}tax\textquotedblright{} is usually traced to Old High German \emph{stiura} \textquotedblleft{}a prop, a support\textquotedblright{}, and so a contribution; EtymDB lists English \emph{steer} as a cousin of the word.
\end{cousinweb}

\begin{grammarlens}[title={What you've built}]
One more word for telling, explaining and describing at B1.
\end{grammarlens}

\subsection*{Your turn}
\begin{itemize}
\raggedright
  \item \emph{Say it:} \emph{die Steuer}
  \item \emph{Say it:} \emph{die Steuer}, once more
  \item \emph{Say it:} a sentence of your own with \emph{die Steuer}
  \item \emph{Recall:} say the German for the fortune, then the German for the payment, then say \emph{die Steuer} again
\end{itemize}

\begin{tcolorbox}[breakable,colback=teal!4,colframe=teal!35!black,title={Before you move on}]
What does die Steuer mean? (\textquotedblleft{}The tax\textquotedblright{}.) Say it once more.
\end{tcolorbox}
\section[der Zins]{der Zins — the interest (on money)}

\label{lesson:GE-C391-zins}



\begin{quote}
Before the new one: say the German for the payment, then the German for the tax.
\end{quote}

\subsection*{You'll want to know: der Zins}
\textbf{der Zins} — \textquotedblleft{}the interest (on money)\textquotedblright{}.

\begin{quote}
Die Bank zahlt nur ein Prozent Zinsen.
\end{quote}

\begin{quote}
\emph{The bank pays only one per cent interest.}
\end{quote}

Mostly plural: \emph{die Zinsen}. In southern Germany and Austria \emph{der Zins} can also mean rent.

\begin{cousinweb}
\emph{Zins} is inherited from Middle High German and Old High German \emph{zins} \textquotedblleft{}tax, tribute, interest\textquotedblright{}, an early borrowing from Latin \emph{cēnsus} \textquotedblleft{}assessment, tax\textquotedblright{} (EtymDB), from \emph{cēnsēre} \textquotedblleft{}to assess\textquotedblright{}.
\end{cousinweb}

\begin{grammarlens}[title={What you've built}]
One more word for telling, explaining and describing at B1.
\end{grammarlens}

\subsection*{Your turn}
\begin{itemize}
\raggedright
  \item \emph{Say it:} \emph{der Zins}
  \item \emph{Say it:} \emph{der Zins}, once more
  \item \emph{Say it:} a sentence of your own with \emph{der Zins}
  \item \emph{Recall:} say the German for the payment, then the German for the tax, then say \emph{der Zins} again
\end{itemize}

\begin{tcolorbox}[breakable,colback=teal!4,colframe=teal!35!black,title={Before you move on}]
What does der Zins mean? (\textquotedblleft{}The interest (on money)\textquotedblright{}.) Say it once more.
\end{tcolorbox}
\section[investieren]{investieren — to invest}

\label{lesson:GE-C391-investieren}



\begin{quote}
Before the new one: say the German for the tax, then the German for the interest.
\end{quote}

\subsection*{You'll want to know: investieren}
\textbf{investieren} — \textquotedblleft{}to invest\textquotedblright{}.

\begin{quote}
Sie hat ihr Geld in eine kleine Firma investiert.
\end{quote}

\begin{quote}
\emph{She invested her money in a small firm.}
\end{quote}

Perfekt: \emph{hat investiert}. \emph{Zeit investieren} — to invest time.

\begin{cousinweb}
\emph{investieren} goes back to Latin \emph{investīre} \textquotedblleft{}to clothe, to cover\textquotedblright{} (Lewis \& Short), built on \emph{vestis} \textquotedblleft{}garment\textquotedblright{}; the money sense is usually said to have come later, through Italian \emph{investire}.
\end{cousinweb}

\begin{grammarlens}[title={What you've built}]
One more word for telling, explaining and describing at B1.
\end{grammarlens}

\subsection*{Your turn}
\begin{itemize}
\raggedright
  \item \emph{Say it:} \emph{investieren}
  \item \emph{Say it:} \emph{investieren}, once more
  \item \emph{Say it:} a sentence of your own with \emph{investieren}
  \item \emph{Recall:} say the German for the tax, then the German for the interest, then say \emph{investieren} again
\end{itemize}

\begin{tcolorbox}[breakable,colback=teal!4,colframe=teal!35!black,title={Before you move on}]
What does investieren mean? (\textquotedblleft{}To invest\textquotedblright{}.) Say it once more.
\end{tcolorbox}
\section[anschaffen]{anschaffen — to acquire, to buy (something lasting)}

\label{lesson:GE-C391-anschaffen}



\begin{quote}
Before the new one: say the German for the interest, then the German for to invest.
\end{quote}

\subsection*{You'll want to know: anschaffen}
\textbf{anschaffen} — \textquotedblleft{}to acquire, to buy (something lasting)\textquotedblright{}.

\begin{quote}
Wir wollen uns eine neue Waschmaschine anschaffen.
\end{quote}

\begin{quote}
\emph{We want to buy a new washing machine.}
\end{quote}

Separable: \emph{wir schaffen an}, \emph{wir haben angeschafft}. Often with \emph{sich}. \emph{Die Anschaffung} — the purchase.

\begin{cousinweb}
\emph{anschaffen} is built from \emph{an} and \emph{schaffen} \textquotedblleft{}to make, to get done\textquotedblright{}, inherited from Old High German \emph{scaffan} and Proto-Germanic \emph{skapjaną} \textquotedblleft{}to create\textquotedblright{} (EtymDB), a cousin of English \emph{shape}.
\end{cousinweb}

\begin{grammarlens}[title={What you've built}]
One more word for telling, explaining and describing at B1.
\end{grammarlens}

\subsection*{Your turn}
\begin{itemize}
\raggedright
  \item \emph{Say it:} \emph{anschaffen}
  \item \emph{Say it:} \emph{anschaffen}, once more
  \item \emph{Say it:} a sentence of your own with \emph{anschaffen}
  \item \emph{Recall:} say the German for the interest, then the German for to invest, then say \emph{anschaffen} again
\end{itemize}

\begin{tcolorbox}[breakable,colback=teal!4,colframe=teal!35!black,title={Before you move on}]
What does anschaffen mean? (\textquotedblleft{}To acquire, to buy (something lasting)\textquotedblright{}.) Say it once more.
\end{tcolorbox}
\section[versichern]{versichern — to insure; to assure}

\label{lesson:GE-C391-versichern}



\begin{quote}
Before the new one: say the German for to invest, then the German for to acquire.
\end{quote}

\subsection*{You'll want to know: versichern}
\textbf{versichern} — \textquotedblleft{}to insure; to assure\textquotedblright{}.

\begin{quote}
Wir haben das Haus gegen Feuer versichert.
\end{quote}

\begin{quote}
\emph{We have insured the house against fire.}
\end{quote}

Inseparable: \emph{er versichert, er hat versichert}. Also to assure: \emph{Ich versichere dir, dass …} \emph{Die Versicherung} — the insurance.

\begin{cousinweb}
\emph{versichern} is built with \emph{ver-} on \emph{sicher} \textquotedblleft{}safe\textquotedblright{}, inherited from Old High German \emph{sihhur}, which goes back to Latin \emph{sēcūrus} \textquotedblleft{}free from care\textquotedblright{} (EtymDB; Lewis \& Short): you make a thing safe.
\end{cousinweb}

\begin{grammarlens}[title={What you've built}]
One more word for telling, explaining and describing at B1.
\end{grammarlens}

\subsection*{Your turn}
\begin{itemize}
\raggedright
  \item \emph{Say it:} \emph{versichern}
  \item \emph{Say it:} \emph{versichern}, once more
  \item \emph{Say it:} a sentence of your own with \emph{versichern}
  \item \emph{Recall:} say the German for to invest, then the German for to acquire, then say \emph{versichern} again
\end{itemize}

\begin{tcolorbox}[breakable,colback=teal!4,colframe=teal!35!black,title={Before you move on}]
What does versichern mean? (\textquotedblleft{}To insure; to assure\textquotedblright{}.) Say it once more.
\end{tcolorbox}
